% Sec. V-E "Synthetic Views from the Reference" body for main.tex. \input
% from main.tex after the Scoring Capability subsection. Tooling:
% ../isaac_coop_sim (export_usd.py, make_scenario.py, run_isaac.py).
The same reference that scores the real runs can also stand in for the
building in a simulator, so that views the testbed did not record, or
recorded too thinly, are generated against the geometry it measured.
The zero-level surface of the signed distance field is exported as a
USD stage in the map frame, with its vertex colours and a collision
surface, and loaded into Isaac Sim \cite{isaacsim}. Each robot is
placed on the trajectory the reference pipeline registered for it
(Sec.~\ref{sec:gt}), so a synthetic run drives through the same doors at
the same times as the real one, and carries the sensors of
Table~\ref{tab:sensors} at the extrinsics of
Table~\ref{tab:calib}: the LiDAR robot a 64-beam and the RGB-D cart a
128-beam ray-traced Ouster model at \SI{10}{\hertz} and 1024 columns,
and the cameras as pinhole models with the calibrated intrinsics
($960\times540$ and $848\times480$ pixels) at \SI{15}{\hertz}. Each robot
is also a body of its measured footprint, so it appears in the other's
scans. Moving actors, people walking waypoint loops at
\SI{1.2}{\metre\per\second} and furniture moved between runs, are
added by the scenario; the real recordings have few of them, since
a run is two robots and the people operating them.

The output is written as a rosbag2 recording with the real topic names,
frames, and transform tree, so it enters the same converter and the same
scoring as a recorded run, and nothing downstream knows the difference.
What differs is the ground truth: in the twin every pose is exact and
every object, moving or not, has a box at every LiDAR stamp with no
annotation, where the real runs label static objects once in the map
frame (Sec.~\ref{sec:gt_score}) and moving ones by hand. Table~\ref{tab:twin}
lists which inputs the twin takes from the testbed and which it
invents. Two uses follow. A sensor stream whose recording fell short of
its rate, as the stereo camera's did at \SI{1.3}{\hertz} in one
cooperative session, is regenerated at its nominal rate along the
recorded trajectory. And the diversity the two sites lack, more actors,
other furniture, other paths, is added without a recording session, for
pre-training a collaborative method that is then fine-tuned and
evaluated on real sequences only. The twin is a geometric one: it
renders the surfaces the map has, with the appearance the scans gave
them, and carries none of the reflectance, dropout, and motion
distortion of the real sensors, so, as with the channel in E5, a
method's score on it is not its score on the testbed, and the gap
between the two is itself a measurement the testbed makes.
\todo{Report the scan-to-map agreement of synthetic against real scans
at the same poses, and the detection score gap, once the runs exist.}

\begin{table}[!t]
\centering
\caption{The Synthetic Twin: Measured Inputs and Invented Ones}
\label{tab:twin}
\begin{tabular}{lll}
\toprule
Element & Source & Status \\
\midrule
Building surface & reference map (TSDF zero level) & measured \\
Appearance & scan intensity / colour per vertex & measured \\
Robot trajectories & registered poses, \SI{10}{\hertz} & measured \\
Sensor mounts & extrinsics (Table~\ref{tab:calib}) & measured \\
Camera intrinsics & calibration & measured \\
LiDAR beam pattern & Ouster model, 64 / 128 beams & modelled \\
Robot bodies & boxes of the measured footprint & modelled \\
Moving actors & walkers, moved furniture & invented \\
Object boxes & every object, every stamp & exact \\
Link state & none; replay (Sec.~\ref{sec:replay}) & from the real run \\
\bottomrule
\end{tabular}
\end{table}
